%% ============================================================
%%  SECCIÓN 2 — Resumen
%% ============================================================

\section{Resumen}
\label{sec:resumen}

En el presente informe de laboratorio se determinaron y contrastaron experimentalmente las curvas características voltaje--corriente ($V$--$I$) de tres elementos con distinta naturaleza de transporte eléctrico: una resistencia lineal de carbón de valor nominal \qty{100}{\ohm}, una lámpara incandescente de filamento de tungsteno y un diodo semiconductor de silicio en polarización directa, empleando un circuito potenciómetro divisor de tensión accionado por reóstato y alimentado por una fuente de corriente continua a \qty{7,5}{\volt}.

En primer lugar, para la resistencia de carbón, en el intervalo de \qtyrange{1,00}{6,49}{\volt}, se verificó de forma rigurosa la validez de la ley de Ohm. La regresión lineal por mínimos cuadrados ordinarios (MCO) de la relación $I = G \cdot V + b$ arrojó una conductancia experimental $G = \qty{9,998 \pm 0,054e-3}{\siemens}$ y una ordenada al origen estadísticamente nula ($b = \qty{0,11 \pm 0,24}{\milli\ampere}$), equivalente a una resistencia experimental de $R = \qty{100,02 \pm 2,55}{\ohm}$ con un coeficiente de determinación $R^2 = 0{,}99985$ y un error relativo porcentual de apenas $E_r = \qty{0,02}{\percent}$ respecto al valor nominal de \qty{100}{\ohm}.

En segundo lugar, para la lámpara incandescente, en el rango de \qtyrange{1,00}{6,43}{\volt}, se demostró un comportamiento fuertemente sublineal y termorresistivo como consecuencia del calentamiento por efecto Joule y el balance radiativo de Stefan-Boltzmann. El ajuste mediante una ley de potencias $I = k V^n$ proporcionó $I = (0{,}146 \pm 0{,}003)\,V^{0{,}573 \pm 0{,}014}$ ($R^2 = 0{,}9988$). El exponente experimental $n = 0{,}573 < 1$ certifica que la resistencia estática $R = V/I$ se incrementa monótonamente en más del \qty{120}{\percent}, transitando desde \qty{6,76 \pm 0,08}{\ohm} a \qty{1,00}{\volt} hasta \qty{14,88 \pm 0,04}{\ohm} a \qty{6,43}{\volt}, alcanzando una potencia disipada máxima de \qty{2,78}{\watt}.

En tercer lugar, para el diodo semiconductor, se constató un comportamiento rectificador superlineal de tipo exponencial. En tensiones inferiores a \qty{0,60}{\volt}, el dispositivo se mantuvo en estado de corte o bloqueo práctico, mientras que para voltajes comprendidos entre \qtyrange{0,61}{0,75}{\volt} la corriente creció abruptamente hasta \qty{0,783}{\ampere}. La resistencia estática colapsó desde \qty{12,2 \pm 0,3}{\ohm} a \qty{0,61}{\volt} hasta \qty{0,96 \pm 0,01}{\ohm} a \qty{0,75}{\volt}. A partir de la extrapolación lineal de la zona de conducción dinámica, se cuantificó una tensión de umbral o codo de conducción de $V_u \approx \qty{0,66 \pm 0,02}{\volt}$, plenamente concordante con la teoría de junturas $p$-$n$ de silicio. Asimismo, se identificó e interpretó rigurosamente la anomalía instrumental del punto experimental $N^\circ 4$ (\qty{0,70}{\volt}; \qty{0,230}{\ampere}).
